\chapter{The CAN Database}
\label{ch:dbc}

A DBC file is the human-readable contract of the bus~\cite{vectordbc}.
This chapter describes the DBC that was brought into alignment with the
firmware of \cref{ch:dev}: the message tree, two representative signal
layouts, the offset convention, the reserved window for new commands, and
the structural correction of ACK/NACK.

\section{Role of the DBC in this project}
Without a DBC, CANalyzer displays identifiers and hex dumps. With a DBC it
displays ``heatsink temperature, \SI{47.3}{\celsius}'' and it can encode a
speed ramp from two physical fields, \emph{final speed in RPM} and
\emph{duration in ms}. The Python GUI of \cref{ch:gui} never packs a
CAN frame itself; it writes system variables, CAPL reads those variables,
and CANalyzer encodes the frame from the DBC. It follows that a DBC which
disagrees with the firmware is not a documentation issue. It is a
functional bug in the operator path.

The file therefore defines, for every message,

\begin{itemize}[leftmargin=1.6em]
  \item the 11-bit identifier;
  \item the DLC in bytes;
  \item each signal's start bit, length, signedness, endianness, factor,
        offset and unit;
  \item a comment that names the matching firmware symbol where one
        exists.
\end{itemize}

\section{Message tree}
\label{sec:tree}

The DBC is organised in three branches, visible in CANalyzer's
symbol explorer (\cref{fig:dbc-tree}).

\begin{figure}[ht]
\centering
\includegraphics[width=0.72\textwidth]{dbc_tree.png}
\caption{DBC message tree as shown by CANalyzer. Diagnostics occupy
\texttt{0x40}--\texttt{0x60}; motor-control commands occupy
\texttt{0x30}--; MCU responses occupy \texttt{0x10} and \texttt{0x11}.}
\label{fig:dbc-tree}
\end{figure}

\paragraph{Diagnostics, \texttt{0x40}--\texttt{0x60}.}
Ten periodic messages, one per register of \cref{tab:diag-map}. Each has
DLC~4 and a single 32-bit little-endian signal. Factors convert the SDK's
internal units into volts, degrees Celsius, milliamperes or RPM so that
the GUI never sees raw ticks.

\paragraph{Motor-control commands, \texttt{0x30}--.}
\texttt{SET\_REG}, \texttt{GET\_REG}, \texttt{EXECUTE\_COMMAND},
\texttt{SET\_RAMP}, the paired and independent current commands, and the
two current ramps. Layouts are message-specific: a bit-mask, a pair of
physical values, a packed PI word.

\paragraph{MCU responses, \texttt{0x10} and \texttt{0x11}.}
\texttt{ACK\_NACK} and \texttt{REG\_VALUE}. These two messages are the
DBC counterpart of \cref{lst:ack}.

\section{Two representative layouts}
\label{sec:layouts}

\subsection{\texttt{EXECUTE\_COMMAND} (identifier \texttt{0x33})}
The payload is a bit-mask spread over two bytes. Each flag is a 1-bit
signal with factor~1, so CAPL and the GUI can write
\texttt{msg.EXE\_START\_MOTOR = 1} without shifting. The bit positions
match \cref{tab:exec-flags}. Unused bits are present in the DBC as
reserved signals; they are transmitted as zero and ignored on reception,
which is what makes a later extension of the mask binary-compatible.

\subsection{\texttt{SET\_RAMP} (identifier \texttt{0x37})}
Two physical signals share the eight-byte payload:

\begin{itemize}[leftmargin=1.6em]
  \item \texttt{FinalSpeed\_execrmp} --- signed, little-endian, factor
        chosen so that the GUI unit is RPM;
  \item \texttt{Duration\_execrmp} --- unsigned, little-endian, unit
        milliseconds.
\end{itemize}

The same pattern---final value plus duration---is reused for
\texttt{SET\_ID\_RAMP} and \texttt{SET\_IQ\_RAMP}, with the final-value
signal scaled in milliamperes rather than RPM. Reusing the pattern is a
deliberate invitation to reuse CAPL: the speed-ramp event handler of
\cref{lst:capl-speed} has a current-ramp twin that differs only in the
message type.

\begin{figure}[ht]
\centering
\includegraphics[width=0.95\textwidth]{dbc_signals.png}
\caption{Signal-level view of a DBC message in CANalyzer. Start bit,
length, factor and offset are the fields that must match the firmware's
little-endian packing.}
\label{fig:dbc-signals}
\end{figure}

\section{Alignment with the $+$\texttt{0x30}/$+$\texttt{0x40} rule}
\label{sec:dbc-offset}

The DBC was rewritten so that a reader who knows the firmware offset rule
of \cref{sec:offset} can predict every identifier, and a reader who knows
the DBC can predict every firmware code. Three ranges result.

\begin{center}
\begin{tabular}{lll}
\toprule
Range & Role & Origin \\
\midrule
\texttt{0x10}--\texttt{0x11} & MCU responses & size-based dispatch, \cref{sec:ack} \\
\texttt{0x30}--\texttt{0x3A} & original commands & firmware library as found \\
\texttt{0x3B}--\texttt{0x3C} & \texttt{SET\_CURRENT\_ID} / \texttt{IQ} & added in this work \\
\texttt{0x3E}--\texttt{0x3F} & $I_d$ / $I_q$ ramps & $I_q$ existing, $I_d$ added \\
\texttt{0x40}-- & diagnostics & \texttt{Send\_Status} \\
\bottomrule
\end{tabular}
\end{center}

\texttt{SET\_CURRENT\_ID} and \texttt{SET\_CURRENT\_IQ} illustrate the one
place where the naive offset was not used. The protocol codes
\texttt{0x0F} and \texttt{0x11} would have mapped to \texttt{0x3F} and
\texttt{0x41}. The first collides with \texttt{SET\_IQ\_RAMP}; the second
sits inside the diagnostic window. Unused identifiers \texttt{0x3B} and
\texttt{0x3C} were therefore assigned in both the DBC and the C
\texttt{switch}, and a comment in each artefact points at the other. The
offset rule remains the default; collisions are resolved by an explicit
exception, never by silent reuse.

\section{Structural correction of ACK/NACK}
\label{sec:dbc-ack}

The original DBC declared two messages with identifiers \texttt{0xF0} and
\texttt{0xFF}. Those values are the \emph{payload} of an acknowledgement,
not CAN identifiers: they do not fit in 11 bits without being interpreted
as a different, extended-id convention, and they never appeared on the
wire as identifiers. CANalyzer therefore never decoded a real ACK, and
any CAPL handler attached to those messages was dead code.

\begin{figure}[ht]
\centering
\includegraphics[width=0.55\textwidth]{ack_before.png}
\caption{Original DBC fragment in which \texttt{0xFF} (NACK) and
\texttt{0xF0} (ACK) are declared as distinct CAN identifiers. Both values
are in fact payload codes of the single message \texttt{0x10}.}
\label{fig:ack-before}
\end{figure}

The repair is a single message \texttt{ACK\_NACK} at identifier
\texttt{0x10}, DLC~1, with one unsigned 8-bit signal
\texttt{AckCode}. A value table maps \texttt{0xF0} $\mapsto$
\texttt{ACK} and \texttt{0xFF} $\mapsto$ \texttt{NACK}. Identifier
\texttt{0x11} remains \texttt{REG\_VALUE}, with two 16-bit signals
\texttt{GET\_K\_NUMERATOR} and \texttt{GET\_K\_DENOMINATOR} that mirror
the packing of \cref{sec:pi-pack} for PI reads, and a 32-bit raw signal
for every other register.

After this change, a CAPL \texttt{on message ACK\_NACK} handler fires on
every write, and a CAPL \texttt{on message REG\_VALUE} handler fires on
every read. The Python GUI polls the corresponding system variables and
never has to look at identifiers.

\section{Keeping the DBC honest}
A DBC drifts as soon as a firmware identifier is added without a matching
edit. The working rule adopted at the end of this project is
mechanical and small:

\begin{itemize}[leftmargin=1.6em]
  \item no firmware identifier ships without a DBC message of the same
        name and DLC;
  \item no DBC signal ships without a comment that names the C symbol or
        the bit it corresponds to;
  \item a PCAN-View trace of a newly added message is attached to the
        laboratory notebook before the GUI is allowed to send it.
\end{itemize}

Those three checks would have caught both the ACK-as-identifier mistake
and the SET\_REG length-byte mistake before either reached an operator.
